\documentclass[11pt]{article}
\usepackage[utf8]{inputenc}
\usepackage[T1]{fontenc}
\usepackage{lmodern}
\usepackage[margin=2.4cm]{geometry}
\usepackage{xcolor}
\usepackage{enumitem}
\usepackage[hidelinks]{hyperref}
\newcommand{\rev}[1]{\par\smallskip\noindent\textit{Reviewer:} #1\par}
\newcommand{\resp}[1]{\par\noindent\textcolor{blue!60!black}{\textbf{Response:}} #1\par\medskip}
\setlength{\parindent}{0pt}
\setlength{\parskip}{0.4em}

\begin{document}
\begin{center}
{\large\textbf{Response to Reviewer}}\\[0.3em]
\textit{Environmental Monitoring and Assessment}\\[0.2em]
Manuscript: ``A Transferable Protocol for Satellite-Based Water Transparency
Monitoring in High-Altitude Oligotrophic Lakes: Harmonized Sentinel-2/Landsat
Retrieval with Rigorous Validation and Uncertainty Quantification in Lake Titicaca
(2013--2024)''
\end{center}

We thank the Reviewer for a careful, constructive and encouraging assessment. We
are grateful for the recognition of the validation hierarchy, the conformal
uncertainty, the use of real in-situ data, and the honest retrievability envelope.
We have addressed every point; in particular, we have substantially strengthened
the framing toward \emph{environmental monitoring systems and operational
assessment}, which was the central request. Below we reproduce each comment and our
response. All changes are incorporated in the revised manuscript.

\hrulefill

\section*{Major comments}

\textbf{1. Strengthen the ``monitoring system'' and ``operational assessment''
framing (critical for EMA).}
\resp{Addressed throughout. (i) We added a dedicated paragraph at the end of the
Introduction that explicitly positions the work as a \emph{remote-sensing
monitoring protocol} complementing and optimizing the periodic field campaigns of
IMARPE, ALT and OEFA, producing synoptic lake-wide transparency maps with quantified
uncertainty \emph{between} campaigns. (ii) We added a new conceptual diagram
(Fig.~1) of the full protocol (field campaigns $\to$ match-ups $\to$ validation
hierarchy $\to$ model + conformal prediction $\to$ map + uncertainty $\to$
operational use). (iii) We added a new Discussion subsection, ``Operational
implications for environmental monitoring,'' detailing how regional authorities
(GORE Puno, ANA, PELT and the binational ALT) can use the maps and uncertainty
intervals for three concrete functions: targeting field sampling by information
value, early warning of trophic-state change (focused on Bah\'ia de Puno), and
transboundary Peru--Bolivia consistency under a single harmonized calibration.}

\textbf{2. Title and Abstract --- make them more impactful and aligned with EMA.}
\resp{We adopted the Reviewer's suggested title, now leading with ``A Transferable
Protocol for Satellite-Based Water Transparency Monitoring in High-Altitude
Oligotrophic Lakes\ldots''. The Abstract was rewritten to lead with the transferable
\emph{monitoring} protocol, to state that it complements in-situ campaigns, and to
make the ``operational envelope'' explicit (delimiting that chlorophyll-$a$,
suspended solids and temperature are not retrievable in this regime), closing on the
prioritization of field sampling between campaigns and the protocol's
transferability.}

\textbf{3. Discussion --- expand operational implications and practical
limitations.}
\resp{The new ``Operational implications'' subsection adds: the operational window
(near-weekly combined Sentinel-2/Landsat cadence vs.\ annual field surveys);
recommended \emph{incremental integration} with existing programmes (each campaign
both validates the current map and extends the training set); and the practical
limitations for operational use --- the dry-season sampling bias, signal saturation
above $\sim$15~m, the positive bias in turbid Bah\'ia de Puno, and an explicit
recommendation: \textbf{local recalibration with field campaigns is required before
applying the protocol to another lake or to an unsampled optical sub-basin}.}

\textbf{4. Add a conceptual diagram of the protocol.}
\resp{Added as Fig.~1 (a new, original schematic). It visualizes the transferable
protocol end-to-end, reinforcing the central contribution exactly as suggested.}

\textbf{5. Figures and visual quality.}
\resp{All figures are provided at 600~dpi. The modelled-transparency map (now
Fig.~8) retains its $Z_{SD}$ colour scale and legend. We additionally signpost the
validation hierarchy (now Fig.~5 and Table~3) in the text as the methodological
core of the protocol, so the reader's attention is drawn to it.}

\section*{Minor comments}

\textbf{English language.}
\resp{We tightened several long sentences (notably the Abstract, the prior-literature
passage in the Introduction, and the Methods). The manuscript has been re-read for
clarity and flow.}

\textbf{References.}
\resp{For the new methodological justification (match-up window) we cite established
operational $Z_{SD}$ references already in the reference list. We are happy to add
any specific EMA-scope references on integrated lake monitoring or uncertainty in
operational products that the Editor may recommend.}

\textbf{Methods --- justify the 30~m buffer and $\pm$10-day window.}
\resp{Done. We now state that the 30~m buffer matches the coarser Landsat pixel and
averages residual geolocation error and sensor noise over several pixels while
remaining in open water, and that the $\pm$10-day window balances scarce cloud-free
availability at 3810~m against the slow temporal change of transparency in this
large, deep lake.}

\textbf{Results --- positive bias in Bah\'ia de Puno ($+1.13$~m) as an operational
limitation.}
\resp{This is now explicitly framed as an operational limitation in the new
Discussion subsection: in the most vulnerable (eutrophying) zone the clear-water
majority pulls predictions upward, so mapped values there must be read together with
their wider uncertainty interval and confirmed by field sampling.}

\textbf{Conclusions.}
\resp{Strengthened to close on the transferable protocol and its operational value:
by delivering maps with calibrated confidence between field campaigns, it turns
intermittent in-situ sampling into a continuous, decision-ready monitoring indicator,
directly transferable to other high-altitude oligotrophic lakes when locally
recalibrated.}

\textbf{Optional --- ``Climate Extremes and Multi-Hazard Events'' call.}
\resp{We note the suggestion. Transparency can serve as an indicator of drought /
extreme-event impact on lake clarity; we have kept this as a brief forward-looking
remark rather than the main focus, to preserve the monitoring-protocol framing,
and we would be glad to develop it further if the Editor considers it appropriate.}

\hrulefill

We believe these revisions materially strengthen the manuscript's alignment with the
scope of \textit{Environmental Monitoring and Assessment} while preserving the
scientific rigour the Reviewer kindly acknowledged. We thank the Reviewer again for
the constructive guidance.

\end{document}
